% Generated by build/sync_qft_soln.py from committed sources only.
% Source repository: https://github.com/Luca-Yucheng-Jin/QFT-soln
% Source commit: 07d7fe95fd5cf5b26d38a16ee535a04a925277b2
\section{$\phi^3$ Theory (PSI QFT II, Tutorial 5)}

\noindent\textit{Question prompts paraphrased from the official
\href{https://psi-online.perimeterinstitute.ca/courses/take/quantum-field-theory-ii-student/pdfs/20776383-module-4-tutorial-2}{Perimeter PSI Online sheet};
the equations and notation follow the source.}

\subsection{The Effective Action for $\phi^3$ Theory}

In six-dimensional Euclidean spacetime, consider
\begin{equation}
    S_\lambda
    =\int d^6x\,
    \left[
        \frac{1}{2}(\partial\phi)^2
        +\frac{m^2}{2}\phi^2
        +\frac{\lambda}{3!}\phi^3
    \right].
\end{equation}
For a bosonic theory, the one-loop effective action is
\begin{equation}
    \Gamma[\phi]
    =S[\phi]
    +\frac{\hbar}{2}\operatorname{Tr}
    \!\left[\ln\!\left(S''[\phi]\right)\right]
    +O(\hbar^2),
\end{equation}
where terms of order $\hbar^n$ correspond to $n$-loop diagrams.

\begin{enumerate}[label=\alph*)]
    \item Calculate the integral kernel associated with the operator
    $S''[\phi]$ in $\phi^3$ theory.
    \begin{solution}
        \begin{equation}
            S''[\phi]_{xy} = \left(-\partial^2 + m^2 + \lambda\phi(x)\right)\delta^d(x-y).
        \end{equation}
    \end{solution}

    \item Factor the free-field contribution out of $S''[\phi]$ and expand
    $\operatorname{Tr}[\ln(S''[\phi])]$ perturbatively in $\lambda$.

    \textit{Hint:} Recall
    \begin{equation}
        \ln(1+x)
        =x-\frac{x^2}{2}+\frac{x^3}{3}-\frac{x^4}{4}
        +\frac{x^5}{5}-\frac{x^6}{6}+\cdots.
    \end{equation}
    \begin{solution}
        Write
        \begin{equation}
            (-\partial^2+m^2)\Delta(x-y) = \delta^d(x-y) \iff (-\partial^2 +m^2)_{xy}\Delta_{yz} = \delta_{xz}.
        \end{equation}
        Then, we have
        \begin{equation}
            S''[\phi]_{xy} = (-\partial^2+m^2)_{xz}\left(\delta_{zy} +  \Delta_{zy} \lambda\phi(z)\right).
        \end{equation}
        This immediately follows
        \begin{equation}
        \begin{aligned}
            \Tr \ln (S''[\phi]) &= \Tr \left(\ln(-\partial^2 +m^2)_{xz} + \ln\left(\delta_{zy}+ \lambda\Delta_{zy}\phi(z)\right)\right)\\
            &=\ln(\det(-\partial^2 +m^2)) + \Tr\left(\lambda\Delta_{zy}\phi(y) - \frac{\lambda^2}{2}\Delta_{zu}\Delta_{uy}\phi(z)\phi(u) + \cdots\right)\\
            &=\ln(\det(-\partial^2 +m^2)) + \lambda\Delta(0)\int d^dx \phi(x) - \frac{\lambda^2}{2}\int d^dx d^dy\phi(x)\Delta^2(x-y)\phi(y) + \cdots.
        \end{aligned}
        \end{equation}
    \end{solution}

    \item Express the complete one-loop effective action $\Gamma[\phi]$ as a
    power series in $\lambda$ through fifth order. Represent the result using
    one-loop amputated Feynman diagrams.
    \begin{solution}
        Up to fifth order in $\lambda$, we have
        \begin{equation}
    \begin{aligned}
        \Gamma[\phi]
        &=\int d^dx\,
        \left[
            \frac{1}{2}(\partial\phi)^2
            +\frac{m^2}{2}\phi^2
            +\frac{\lambda}{3!}\phi^3
            +\frac{\hbar\lambda}{2}\Delta(0)\phi(x)
        \right]\\
        &\quad
        -\frac{\hbar\lambda^2}{4}
        \int d^dx\,d^dy\,
        \phi(x)\Delta^2(x-y)\phi(y)\\
        &\quad
        +\frac{\hbar\lambda^3}{6}
        \int d^dx\,d^dy\,d^dz\Delta(x-y)\Delta(y-z)\Delta(z-x)
        \phi(x)\phi(y)\phi(z)\\
        &\quad
        -\frac{\hbar\lambda^4}{8}
        \int d^dx\,d^dy\,d^dz\,d^dw
        \Delta(x-y)\Delta(y-z)\Delta(z-w)\Delta(w-x)
        \phi(x)\phi(y)\phi(z)\phi(w)\\
        &\quad
        +\frac{\hbar\lambda^5}{10}
        \int d^dx\,d^dy\,d^dz\,d^dw\,d^du
        \Delta(x-y)\Delta(y-z)\Delta(z-w)
        \Delta(w-u)\Delta(u-x)\\
        &\qquad \qquad \times
        \phi(x)\phi(y)\phi(z)\phi(w)\phi(u)\\
        &\quad
        +\mathcal{O}(\lambda^6)+\mathcal{O}(\hbar^2).
    \end{aligned}
    \end{equation}
        where we have ignore overall constants. These corresponds to the diagrams
            \begin{center}
    \begin{tikzpicture}[
        baseline=(current bounding box.center),
        tutorialFiveScalar/.style={
            draw=black,
            line width=0.8pt
        },
        tutorialFiveVertex/.style={
            circle,
            fill=black,
            draw=black,
            minimum size=3.2pt,
            inner sep=0pt
        },
        every node/.style={font=\scriptsize}
    ]

        % One-point tadpole
        \begin{scope}
            \begin{feynman}
                \vertex (e1) at (0.15,1.8) {$\phi(x)$};
                \vertex[tutorialFiveVertex] (v1) at (1.15,1.8) {};

                \diagram*{
                    (e1) -- [tutorialFiveScalar] (v1)
                };

                \draw[tutorialFiveScalar]
                    (v1) ++(0.45,0) circle (0.45);
            \end{feynman}

        \end{scope}

        % Two-point bubble
        \begin{scope}
            \begin{feynman}
                \vertex (e2x) at (3.55,1.8) {$\phi(x)$};
                \vertex[tutorialFiveVertex] (v2x) at (4.25,1.8) {};
                \vertex[tutorialFiveVertex] (v2y) at (5.35,1.8) {};
                \vertex (e2y) at (6.05,1.8) {$\phi(y)$};

                \diagram*{
                    (e2x) -- [tutorialFiveScalar] (v2x),
                    (v2y) -- [tutorialFiveScalar] (e2y)
                };

                \draw[tutorialFiveScalar]
                    (v2x) to[bend left=55] (v2y);

                \draw[tutorialFiveScalar]
                    (v2x) to[bend right=55] (v2y);
            \end{feynman}

        \end{scope}

        % Three-point triangle
        \begin{scope}
            \begin{feynman}
                \vertex[tutorialFiveVertex] (v3x) at (9.00,2.30) {};
                \vertex[tutorialFiveVertex] (v3y) at (8.50,1.43) {};
                \vertex[tutorialFiveVertex] (v3z) at (9.50,1.43) {};

                \vertex (e3x) at (9.00,3.00) {$\phi(x)$};
                \vertex (e3y) at (7.85,1.05) {$\phi(y)$};
                \vertex (e3z) at (10.15,1.05) {$\phi(z)$};

                \diagram*{
                    (v3x) -- [tutorialFiveScalar] (v3y)
                          -- [tutorialFiveScalar] (v3z)
                          -- [tutorialFiveScalar] (v3x),
                    (e3x) -- [tutorialFiveScalar] (v3x),
                    (e3y) -- [tutorialFiveScalar] (v3y),
                    (v3z) -- [tutorialFiveScalar] (e3z)
                };
            \end{feynman}

        \end{scope}

        % Four-point box
        \begin{scope}
            \begin{feynman}
                \vertex[tutorialFiveVertex] (v4x) at (2.70,-1.05) {};
                \vertex[tutorialFiveVertex] (v4y) at (3.50,-1.05) {};
                \vertex[tutorialFiveVertex] (v4z) at (3.50,-1.85) {};
                \vertex[tutorialFiveVertex] (v4w) at (2.70,-1.85) {};

                \vertex (e4x) at (2.05,-0.40) {$\phi(x)$};
                \vertex (e4y) at (4.15,-0.40) {$\phi(y)$};
                \vertex (e4z) at (4.15,-2.50) {$\phi(z)$};
                \vertex (e4w) at (2.05,-2.50) {$\phi(w)$};

                \diagram*{
                    (v4x) -- [tutorialFiveScalar] (v4y)
                          -- [tutorialFiveScalar] (v4z)
                          -- [tutorialFiveScalar] (v4w)
                          -- [tutorialFiveScalar] (v4x),
                    (e4x) -- [tutorialFiveScalar] (v4x),
                    (e4y) -- [tutorialFiveScalar] (v4y),
                    (v4z) -- [tutorialFiveScalar] (e4z),
                    (v4w) -- [tutorialFiveScalar] (e4w)
                };
            \end{feynman}

        \end{scope}

        % Five-point pentagon
        \begin{scope}
            \begin{feynman}
                \vertex[tutorialFiveVertex] (v5x) at (8.00,-0.85) {};
                \vertex[tutorialFiveVertex] (v5y) at (8.57,-1.27) {};
                \vertex[tutorialFiveVertex] (v5z) at (8.35,-1.95) {};
                \vertex[tutorialFiveVertex] (v5w) at (7.65,-1.95) {};
                \vertex[tutorialFiveVertex] (v5u) at (7.43,-1.27) {};

                \vertex (e5x) at (8.00,-0.05) {$\phi(x)$};
                \vertex (e5y) at (9.35,-0.95) {$\phi(y)$};
                \vertex (e5z) at (8.82,-2.68) {$\phi(z)$};
                \vertex (e5w) at (7.18,-2.68) {$\phi(w)$};
                \vertex (e5u) at (6.65,-0.95) {$\phi(u)$};

                \diagram*{
                    (v5x) -- [tutorialFiveScalar] (v5y)
                          -- [tutorialFiveScalar] (v5z)
                          -- [tutorialFiveScalar] (v5w)
                          -- [tutorialFiveScalar] (v5u)
                          -- [tutorialFiveScalar] (v5x),
                    (e5x) -- [tutorialFiveScalar] (v5x),
                    (v5y) -- [tutorialFiveScalar] (e5y),
                    (v5z) -- [tutorialFiveScalar] (e5z),
                    (v5w) -- [tutorialFiveScalar] (e5w),
                    (e5u) -- [tutorialFiveScalar] (v5u)
                };
            \end{feynman}

        \end{scope}

    \end{tikzpicture}
    \end{center}
    \end{solution}

    \item Explain how the coefficients in the expansion from the preceding part
    are related to the symmetries of the corresponding graphs.

    \item Use the one-loop effective action to show that quantum effects in
    $\phi^3$ theory generate a vacuum expectation value for $\phi(x)$.

    \textit{Hint:} The quantum vacuum field $\varphi(x)$ is the configuration
    that extremizes the quantum effective action.
    \begin{solution}
        Since there is now a term of order 1 in $\phi$, there is a constant term in the effective equation of motion, $\Gamma'[\varphi] = 0$. Therefore, $\varphi$ must be nonzero, and by the definition of $\varphi$, the VEV of $\phi^3$ theory is nonzero.
    \end{solution}
\end{enumerate}

\subsection{Renormalization}

\begin{enumerate}[label=\alph*)]
    \item Redefine the field so that it creates a normalized one-particle state
    with no vacuum overlap:
    \begin{equation}
        \langle\Omega|\phi(0)|\Omega\rangle=0,
        \qquad
        \langle\mathbf k|\phi(0)|\Omega\rangle=1.
    \end{equation}
    Show that the Lagrangian can then be written as
    \begin{equation}
        \mathcal L
        =\frac{1}{2}A\,\partial_\mu\phi\,\partial^\mu\phi
        +B\phi
        +\frac{1}{2}C\phi^2
        +\frac{1}{3!}D\phi^3.
    \end{equation}
    \begin{solution}
        In general, the bare field would give
        \begin{equation}
            \langle\Omega|\phi_0(0)|\Omega\rangle=v,
        \qquad
        \langle\mathbf k|\phi_0(0)|\Omega\rangle=\sqrt{Z}.
        \end{equation}
        Define the renormalised field as
        \begin{equation}
            \phi(x) = \frac{1}{\sqrt{Z}}\phi_0(x) -v,
        \end{equation}
        so that
        \begin{equation}
            \langle\Omega|\phi(0)|\Omega\rangle=0,
            \qquad
            \langle\mathbf k|\phi(0)|\Omega\rangle=1.
        \end{equation}
        Plugging this into the Lagrangian, we find
        \begin{equation}
                    \mathcal L
        =\frac{1}{2}A\,\partial_\mu\phi\,\partial^\mu\phi
        +B\phi
        +\frac{1}{2}C\phi^2
        +\frac{1}{3!}D\phi^3
        \end{equation}
        after some redefinition.
    \end{solution}

    \item The factor $A$ ensures the equal-time canonical commutator
    $[\phi(\mathbf x),\Pi(\mathbf y)]=i\hbar\delta(\mathbf x-\mathbf y)$.
    Use the six-dimensional K\"all\'en--Lehmann representation
    \begin{align}
        \langle\Omega|\phi(x)\phi(y)|\Omega\rangle
        &={}
        \int\frac{d^5\mathbf k}{(2\pi)^5}\,
        \frac{e^{ik\cdot(x-y)}}{2E_{\mathbf k}}
        \left|\langle\mathbf k|\phi(0)|\Omega\rangle\right|^2
        \notag\\
        &\quad+
        \int_{4m^2}^{\infty}ds\,
        \int\frac{d^5\mathbf k}{(2\pi)^5}\,
        \frac{e^{ik\cdot(x-y)}}{2E_{\mathbf k}}\rho(s)
    \end{align}
    to find the relation between $A$ and $\rho(s)$.
    \begin{solution}
        Consider
        \begin{equation}
            \begin{aligned}
                \partial_{y_0}\bra{\Omega}[\phi(x),\phi(y)]\ket{\Omega}\bigg|_{x_0 = y_0} &= \frac{i}{2}\int\frac{d^5\mathbf k}{(2\pi)^5}(e^{-i\mathbf{k}\cdot (\mathbf{x}-\mathbf{y})} + e^{i\mathbf{k}\cdot (\mathbf{x}-\mathbf{y})})\\
                &\qquad+ \frac{i}{2}\int_{4m^2}^{\infty}ds\,
        \int\frac{d^5\mathbf k}{(2\pi)^5}\,
        (e^{-i\mathbf{k}\cdot (\mathbf{x}-\mathbf{y})} + e^{i\mathbf{k}\cdot (\mathbf{x}-\mathbf{y})})\rho(s)\\
        &= i\delta^5(\mathbf{x}- \mathbf{y}) + i\int_{4m^2}^{\infty}ds\,\rho(s)\delta^5(\mathbf{x}- \mathbf{y}).
            \end{aligned}
        \end{equation}
        Meanwhile, we have
        \begin{equation}
            \Pi(x) = Z \partial_0\phi(x).
        \end{equation}
        Therefore, we have
        \begin{equation}
            \partial_{y_0}\bra{\Omega}[\phi(x),\phi(y)]\ket{\Omega}\bigg|_{x_0 = y_0} = \frac{1}{Z}\partial_{y_0}\bra{\Omega}[\phi(\mathbf x),\Pi(\mathbf y)]\ket{\Omega} \implies 1 = Z\left(1+\int_{4m^2}^{\infty}ds\,\rho(s)\right).
        \end{equation}
    \end{solution}

    \item Impose four renormalization conditions to determine $A,B,C,D$. Two
    are the normalization conditions above; a third sets the physical mass to
    zero. For the fourth, require
    \begin{equation}
        \Gamma(p_1,p_2,p_3)=\lambda_R
        \qquad\text{when}\qquad
        (p_1+p_2)^2=(p_1+p_3)^2=\mu^2.
    \end{equation}
    With a sharp momentum cutoff $\Lambda$, determine $A,B,C$ in terms of
    $\lambda_R$ and $\Lambda$ to the order requested, and determine $D$ through
    order $\lambda_R^2$.
    \begin{solution}
        First, let us rewrite the Lagrangian as
        \begin{equation}
            \mathcal{L} = \frac{1}{2}(\partial\phi)^2 + \frac{1}{2}m^2 \phi^2 + \frac{\lambda}{3!}\phi^3 + \frac{\delta_A}{2}(\partial\phi)^2 + B\phi + \frac12 \delta_C \phi^2 + \frac{\delta_D}{3!}\phi^3,
        \end{equation}
        where $\delta_A = A-1$, $\delta_C = C - m^2$, and $\delta_D = D - \lambda$. Let us start with $B$, at one loop order, the only contributing diagram is
\begin{equation}
\begin{aligned}
&\vcenter{\hbox{%
\begin{tikzpicture}[
    baseline=(current bounding box.center),
    phiThreeLine/.style={draw=black,line width=0.85pt},
    phiThreeVertex/.style={
        circle,
        fill=black,
        draw=black,
        minimum size=3.2pt,
        inner sep=0pt
    },
    every node/.style={font=\scriptsize}
]
    \begin{feynman}
        \vertex (ext) at (-1.05,0) {};
        \vertex[phiThreeVertex] (v) at (0,0) {};

        \diagram*{
            (ext) -- [phiThreeLine] (v)
        };

        \draw[phiThreeLine]
            (v) ++(0.42,0) circle (0.42);
    \end{feynman}
\end{tikzpicture}}}
+
\vcenter{\hbox{%
\begin{tikzpicture}[
    baseline=(current bounding box.center),
    phiThreeLine/.style={draw=black,line width=0.85pt},
    every node/.style={font=\scriptsize}
]
    \begin{feynman}
        \vertex (ext) at (-1.05,0) {};
        \vertex (ct) at (0,0);

        \diagram*{
            (ext) -- [phiThreeLine] (ct)
        };

        \def\ctlen{0.13}

        \draw[line width=0.9pt]
            ($(ct)+(-\ctlen,-\ctlen)$)
            --
            ($(ct)+(\ctlen,\ctlen)$);

        \draw[line width=0.9pt]
            ($(ct)+(-\ctlen,\ctlen)$)
            --
            ($(ct)+(\ctlen,-\ctlen)$);
    \end{feynman}
\end{tikzpicture}}}
&=
-\frac{\lambda}{2}
\int\frac{d^6k}{(2\pi)^6}
\frac{1}{k^2+m^2}
-B.
\end{aligned}
\end{equation}
        Since the vacuum expectation value must vanishes, the above must vanishes, and hence we have
        \begin{equation}
            B = -\frac{\lambda}{128\pi^3}\left[
                \frac{\Lambda^4}{4}
                -\frac{m^2\Lambda^2}{2}
                +\frac{m^4}{2}
                \ln\left(1+\frac{\Lambda^2}{m^2}\right)
                \right]
        \end{equation}
        at one-loop order.

        Now, we turn to the mass renormalisation factor. At 1-loop order, we have
\begin{equation}
\begin{aligned}
-M_2(p^2)&=
\vcenter{\hbox{%
\begin{tikzpicture}[
    baseline=(current bounding box.center),
    phiThreeLine/.style={draw=black,line width=0.85pt},
    phiThreeVertex/.style={
        circle,
        fill=black,
        draw=black,
        minimum size=3.2pt,
        inner sep=0pt
    },
    every node/.style={font=\scriptsize}
]
    \begin{feynman}
        \vertex (left) at (-1.50,0) {};
        \vertex[phiThreeVertex] (vL) at (-0.55,0) {};
        \vertex[phiThreeVertex] (vR) at (0.55,0) {};
        \vertex (right) at (1.50,0) {};

        \diagram*{
            (left) -- [phiThreeLine] (vL),
            (vR) -- [phiThreeLine] (right)
        };

        \draw[phiThreeLine]
            (vL) to[bend left=55] (vR);

        \draw[phiThreeLine]
            (vL) to[bend right=55] (vR);
    \end{feynman}
\end{tikzpicture}}}
+
\vcenter{\hbox{%
\begin{tikzpicture}[
    baseline=(current bounding box.center),
    phiThreeLine/.style={draw=black,line width=0.85pt},
    every node/.style={font=\scriptsize}
]
    \begin{feynman}
        \vertex (left) at (-1.25,0) {};
        \vertex (ct) at (0,0);
        \vertex (right) at (1.25,0) {};

        \diagram*{
            (left) -- [phiThreeLine] (ct),
            (ct) -- [phiThreeLine] (right)
        };

        \def\ctlen{0.13}

        \draw[line width=0.9pt]
            ($(ct)+(-\ctlen,-\ctlen)$)
            --
            ($(ct)+(\ctlen,\ctlen)$);

        \draw[line width=0.9pt]
            ($(ct)+(-\ctlen,\ctlen)$)
            --
            ($(ct)+(\ctlen,-\ctlen)$);
    \end{feynman}
\end{tikzpicture}}}\\
&=
\frac{(-\lambda)^2}{2}
\int\frac{d^6k}{(2\pi)^6}
\frac{1}{
    (k^2+m^2)\bigl((k+p)^2+m^2\bigr)
}
-\delta_C-\delta_Ap^2.
\end{aligned}
\end{equation}
        The renormalisation condition requires
        \begin{equation}
            m^2 + M_2(0) = 0 , \quad \lim_{p^2 \to 0} \frac{p^2}{p^2  +m^2 + M_2(p^2)} = 1 \implies M_2(0) = -m^2, \quad M_2'(0) =0.
        \end{equation}
        Therefore, we have
        \begin{equation}
    C
    =
    \frac{\lambda^2}{2}I_2(0)
    =
    \frac{\lambda^2}{256\pi^3}
    \left[
        \Lambda^2
        +\frac{m^2\Lambda^2}{\Lambda^2+m^2}
        -2m^2
        \ln\left(1+\frac{\Lambda^2}{m^2}\right)
    \right],
\end{equation}
\begin{equation}
    A =1+\frac{\lambda^2}{2}I'(0)
    =
    1
    -\frac{\lambda^2}{768\pi^3}
    \left[
        \ln\left(1+\frac{\Lambda^2}{m^2}\right)
        +\frac{2m^2}{\Lambda^2+m^2}
        -\frac{m^4}{2(\Lambda^2+m^2)^2}
        -\frac{3}{2}
    \right].
\end{equation}
Since the first loop correction to the three-point function is of order $\lambda^3$. At order $\lambda^2$, we merely have $D = \lambda_R$.


    \end{solution}
\end{enumerate}
